% ==============================================================================
% THEOREM ENVIRONMENTS
% ==============================================================================
% Provides three theorem-creation commands wrapping keytheorems:
%   \NewNumberedTheorem  — numbered theorem
%   \NewStarredTheorem   — unnumbered theorem
%   \NewDocumentTheorem  — both numbered and starred variants at once
%
% When the framedtheorems package option is active, also loads mdframed,
% tikz, and pgf, and provides \framesetup and \MakeFramedTheorem for
% wrapping theorems in mdframed boxes.
% ==============================================================================

% -----------------------------------------------------------------------
% \NewNumberedTheorem{<cmd>}{<display name>}[<options>]
%   Create a numbered theorem environment via \newkeytheorem.
%   #1 - environment command to define (e.g. \theorem)
%   #2 - human-readable display name (e.g. Theorem)
%   #3 - (optional) key-value options passed to \newkeytheorem
\NewDocumentCommand\NewNumberedTheorem{mmO{}}{
  \newkeytheorem{#1}[name={#2},#3]
}

% -----------------------------------------------------------------------
% \NewStarredTheorem{<cmd>}{<display name>}[<options>]
%   Create an unnumbered (starred) theorem environment via \newkeytheorem.
%   #1 - environment command to define
%   #2 - human-readable display name
%   #3 - (optional) key-value options passed to \newkeytheorem
\NewDocumentCommand\NewStarredTheorem{mmO{}}{
  \newkeytheorem{#1}[name={#2},numbered=false,#3]
}

% -----------------------------------------------------------------------
% \NewDocumentTheorem{<cmd>}{<display name>}[<opts>][<num opts>][<star opts>]
%   Create both numbered (#1) and starred (#1*) theorem environments in one
%   call. Separate option sets can be supplied for each variant.
%   #1 - base environment command (e.g. \theorem; starred becomes \theorem*)
%   #2 - human-readable display name
%   #3 - (optional) key-value options common to both variants
%   #4 - (optional) key-value options for the numbered variant only
%   #5 - (optional) key-value options for the starred variant only
\NewDocumentCommand\NewDocumentTheorem{mmO{}d[]d[]}{
  \IfNoValueTF{#4}
    {\newkeytheorem{#1}[name={#2},#3]}
    {\newkeytheorem{#1}[name={#2},#4,#3]}
  \IfNoValueTF{#5}
    {\newkeytheorem{#1*}[name={#2},numbered=false,#3]}
    {\newkeytheorem{#1*}[name={#2},numbered=false,#5,#3]}
}

% ==============================================================================
% FRAMED THEOREMS (conditional on framedtheorems option)
% ==============================================================================

\if@tssuite@framedtheorems
\ExplSyntaxOn

% -----------------------------------------------------------------------
% Redefine memoize stubs (forward-declared in tssuite.sty)
% -----------------------------------------------------------------------

% When memoize is loaded, disable it during mdframed rendering to avoid
% caching issues with tikz pictures inside framed theorems.
\IfPackageLoadedTF{memoize}{
  \cs_set_protected:Npn\tssuite_memoize_disable:{\mmzset{disable}}
  \cs_set_protected:Npn\tssuite_memoize_enable:{\mmzset{enable}}
}{
  \cs_set_protected:Npn\tssuite_memoize_disable:{}
  \cs_set_protected:Npn\tssuite_memoize_enable:{}
}

% -----------------------------------------------------------------------
% Internal functions
% -----------------------------------------------------------------------

% \tssuite_framesetup:n <title>
%   Configure mdframed appearance. When <title> is blank, a simple frame
%   with no title is set up. When <title> is provided, a tikz node is
%   attached as the frametitle.
%   #1 - frame title text (blank = no title)
\cs_new_protected:Npn\tssuite_framesetup:n#1{
  \tl_if_blank:nTF{#1}
    {
      \mdfsetup{
        innertopmargin=-3pt,
        linecolor=black,
        linewidth=1pt,
        topline=true,
        frametitleaboveskip=\dimexpr-\ht\strutbox\relax
      }
    }
    {
      \mdfsetup{
        frametitle={
          \tikz[baseline=(current~bounding~box.east),outer~sep=0pt]
            \node[anchor=east,rectangle,fill=white]{\strut #1};
        }
      }
      \mdfsetup{
        innertopmargin=-10pt,
        linecolor=black,
        linewidth=1pt,
        topline=true,
        frametitleaboveskip=\dimexpr-\ht\strutbox\relax
      }
    }
}

% -----------------------------------------------------------------------
% User commands
% -----------------------------------------------------------------------

% \framesetup{<title>}
%   User-level wrapper for \tssuite_framesetup:n.
%   #1 - frame title text (blank = no title)
\NewDocumentCommand\framesetup{m}{
  \tssuite_framesetup:n{#1}
}

% -----------------------------------------------------------------------
% \MakeFramedTheorem{<env>}{<theorem cmd>}
%   Create a new environment <env> that wraps an existing theorem
%   (defined via \NewNumberedTheorem or similar) inside an mdframed box.
%
%   Usage:
%     \begin{<env>}[<frame title>][<theorem options>]
%       ...
%     \end{<env>}
%
%   #1 - new environment name (e.g. framedtheorem)
%   #2 - existing theorem command (e.g. \theorem)
\NewDocumentCommand\MakeFramedTheorem{mm}{
  \NewDocumentEnvironment{#1}{oO{}}
    {
      \tssuite_memoize_disable:
      \IfNoValueTF{##1}
        {
          \framesetup{}
          \begin{mdframed}
          \begin{#2}[##2]
        }
        {
          \framesetup{##1}
          \begin{mdframed}
          \begin{#2}[note=##1,##2]
        }
    }
    {
      \end{#2}
      \end{mdframed}
      \tssuite_memoize_enable:
    }
}

\ExplSyntaxOff
\fi
